\section{Discussion and Limitations}
\label{sec:discussion}

\authornote{\textbf{What the results mean.} One paragraph. The isotropic posterior
approximation, not the prior, explains the step-count degradation, because the same
prior with a frequency-resolved covariance is step-robust. Practical consequence:
report PSNR-vs-$T$ curves for any diffusion MRI sampler, and prefer covariances matched
to the forward operator's basis. Tie back to \citet{peng2024optimal}: their argument
was for a better covariance in general; ours is for the \emph{operator-matched} basis
in MRI.}

\authornote{\textbf{Limitations (be specific; reviewers reward this).}
(1) $\Psi$ is estimated from signal-free regions because public raw datasets
(fastMRI, M4Raw) document no noise prescans; mis-estimation is studied only through
$\sigma_y$ scaling (rows h--i).
(2) The Gaussian reference model sets the \emph{shape} of $\bP_t$. Diffusion posteriors
are not Gaussian, and the frequency gate (if kept) is empirical.
(3) Undersampling is retrospective. Prospective acquisition and non-Cartesian
trajectories are untested; the latter would need a NUFFT and lose the exact
diagonalisation.
(4) No reader study. Box metrics from fastMRI+ come from a single reader and are not
exhaustive \citep{zhao2022fastmriplus}, so they do not establish diagnostic benefit.
(5) M4Raw uses a 3\,T-trained brain prior (field-strength shift).
(6) PSNR/SSIM on magnitude images is biased at low SNR \citep{constantinides1997snr}.}

\authornote{\textbf{Future work} (two sentences): downstream quantitative endpoints
(cartilage segmentation and T2 on SKM-TEA, \citealp{desai2021skmtea}; K2S,
\citealp{tolpadi2023k2s}), where K2S found pixel metrics uncorrelated with segmentation
quality; corrupted-sample robustness (spikes, motion lines,
\citealp{campbellwashburn2016spike,levac2024motion}), which residual refinement could
address.}

\paragraph{Conclusion.}
\authornote{Three sentences: the problem, the operator-matched fix, and the headline
numbers (E1, E2, E3) with CIs.}
